\section{Resolution models and homotopical constancy}

These results construct and characterize the resolution presentation $\sSpace^{\dag\vee}_{\CSS}$. Its $\CSS$ subscript is the manuscript notation for the homotopically constant resolution model; it is distinct from the ordinary complete-Segal-space model used later.

\begin{resultblock}
\begin{proposition}[External tensor and evaluation Quillen adjunctions]
\label{catalogue:A1}
For $A\in\sSet^{\dag}$, $K\in\sSet$, and $W\in\sSpace^{\dag\vee}$, the external tensor and mapping simplicial set satisfy a natural bijection $\Hom(A\boxtimes K,W)\cong\Hom(K,\Map(A,W))$.
For every $n\geq0$, specializing to $K=\Delta^n$ gives $(-)\boxtimes\Delta^n\dashv\ev_n$.
With the Part I dagger Joyal model on the coefficient category and the Reedy model on simplicial diagrams, each of these evaluation adjunctions is Quillen.
The constant-diagram functor is naturally isomorphic to $(-)\boxtimes\Delta^0$.
\end{proposition}

\Source{DaggerModels.PartII.externalTensorAdjunction;
DaggerModels.PartII.evaluationAdjunction;
DaggerModels.PartII.evaluationAdjunction_isQuillen;
DaggerModels.PartII.constantDiagramIsoExternalTensorZero}{lean/DaggerModelsPartII/ResolutionAdjunction.lean;
lean/DaggerModelsPartII/ReedyModelComparison.lean}
\CatalogueNote{Scope}{arbitrary u}
\CatalogueNote{Manuscript context}{cs.lem.adjunction, cs.cor.evk, cs.lem.evaluation-quillen}
\CatalogueNote{Construction and reuse}{Project-specific external tensor and adjunction, combined with the project Reedy model and Part I coefficient model.}
\end{resultblock}

\begin{resultblock}
\begin{proposition}[Left proper combinatorial Reedy model]
\label{catalogue:A2}
The category $\Fun(\prism^{\mathrm{op}},\sSet^{\dag})$ has a Reedy model structure whose weak equivalences are the levelwise weak equivalences of the Part I dagger Joyal model.
Pushout along a Reedy cofibration preserves those weak equivalences.
The diagram category is locally presentable, and the two weak factorization systems have small generating families obtained from the Part I free dagger boundary and Smith generators.
These families admit the small object argument, so the Reedy structure is combinatorial.
\end{proposition}

\Source{DaggerModels.PartII.resolutionReedyModelCategory;
DaggerModels.PartII.resolutionReedy_leftProper;
DaggerModels.PartII.resolutionReedy_combinatorial}{lean/DaggerModelsPartII/ReedyProperties.lean;
lean/DaggerModelsPartII/ReedyModelComparison.lean;
lean/DaggerModelsPartII/ReedyCombinatorial.lean}
\CatalogueNote{Scope}{arbitrary u}
\CatalogueNote{Manuscript context}{cs.lem.reedy}
\CatalogueNote{Construction and reuse}{Specializes the project's generic simplicial Reedy construction; left properness and explicit small generators are assembled for the Part I dagger Joyal coefficients.}
\end{resultblock}

\begin{resultblock}
\begin{proposition}[Collapse maps detect all simplicial operators]
\label{catalogue:A3}
Let $G_{\dag}$ consist of the free dagger boundaries and free dagger simplices, and let $S_h$ consist of their external tensors with the collapse maps $\Delta^k\to\Delta^0$ for $k\geq1$.
For a morphism property containing identities and satisfying two-out-of-three, the collapse tests for a fixed coefficient object are equivalent to tests on all ordinary simplex operators.
Consequently, for a Reedy fibrant diagram $W$, locality with respect to the literal family $S_h$ is equivalent to the corresponding derived mapping-space tests for every simplex operator.
\end{proposition}

\Source{DaggerModels.PartII.tensorCollapseTests_iff_allOperators;
DaggerModels.PartII.resolutionLocalObject_iff_allOperators}{lean/DaggerModelsPartII/ResolutionGenerators.lean;
lean/DaggerModelsPartII/ResolutionLocality.lean}
\CatalogueNote{Scope}{arbitrary u}
\CatalogueNote{Manuscript context}{cs.not.Gh, cs.thm.resolution}
\CatalogueNote{Construction and reuse}{Project-specific reduction for the manuscript's literal localization maps; two-out-of-three is the abstract input.}
\end{resultblock}

\begin{resultblock}
\begin{proposition}[Free dagger generators detect weak equivalences]
\label{catalogue:A4}
Suppose $X$ and $Y$ are fibrant in the Part I dagger Joyal model and $f:X\to Y$ is a morphism.
If, for every $n\geq0$, the derived mapping simplicial sets from both the free dagger boundary of $\Delta^n$ and the free dagger simplex $\Delta^n$ send $f$ to a Kan weak equivalence, then $f$ is a dagger Joyal weak equivalence.
The proof uses the actual Part I generating cofibrations and their cofibrant endpoints;
it supplies the converse direction of the resolution-locality characterization.
\end{proposition}

\Source{DaggerModels.PartII.weak_of_resolutionTestObjects}{lean/DaggerModelsPartII/ResolutionLocalCharacterization.lean}
\CatalogueNote{Scope}{arbitrary u}
\CatalogueNote{Manuscript context}{cs.thm.resolution}
\CatalogueNote{Construction and reuse}{Project-specific specialization of a generic derived-mapping detection argument to the actual Part I free dagger boundary and simplex generators.}
\end{resultblock}

\begin{resultblock}
\begin{proposition}[Local and fibrant objects are homotopically constant]
\label{catalogue:A5}
For a simplicial dagger diagram $W$, locality with respect to $S_h$ is equivalent to Reedy fibrancy together with a levelwise Part I dagger Joyal weak equivalence $\ell_W:cW_0\to W$.
After constructing the exact localized model, its fibrant objects satisfy precisely this same condition.
Thus the characterization applies to the specified localization itself, without an additional hypothesis that $W$ is local.
The levelwise counit condition is also equivalent to requiring every simplicial structure map of $W$ to be a weak equivalence.
\end{proposition}

\Source{DaggerModels.PartII.resolutionLocalObject_iff_homotopicallyConstant;
DaggerModels.PartII.resolutionLocalized_isFibrant_iff_local;
DaggerModels.PartII.resolutionLocalized_isFibrant_iff;
DaggerModels.PartII.resolutionCounit_levelwise_iff}{lean/DaggerModelsPartII/ResolutionLocalCharacterization.lean;
lean/DaggerModelsPartII/ResolutionModel.lean;
lean/DaggerModelsPartII/ResolutionAdjunction.lean}
\CatalogueNote{Scope}{arbitrary u}
\CatalogueNote{Manuscript context}{cs.thm.resolution}
\CatalogueNote{Construction and reuse}{Project-specific detection using the literal $S_h$ family, then applied to the constructed local model; the final structure-map equivalence uses two-out-of-three.}
\end{resultblock}

\begin{resultblock}
\begin{proposition}[Local equivalences between local objects are detected at zero]
\label{catalogue:A6}
If $W$ and $V$ are $S_h$-local Reedy fibrant diagrams and $f:W\to V$, then $f$ is an $S_h$-local equivalence if and only if its zero component $f_0:W_0\to V_0$ is a weak equivalence in the Part I dagger Joyal model.
More precisely, local equivalence first agrees with Reedy weak equivalence between local objects, and homotopical constancy then reduces the levelwise condition to degree zero.
The statement does not assert this zero-degree criterion for arbitrary diagrams without local replacement.
\end{proposition}

\Source{DaggerModels.PartII.resolutionLocalEquivalence_iff_reedyWeak_of_local;
DaggerModels.PartII.resolutionLocalEquivalence_iff_evalZero_of_local}{lean/DaggerModelsPartII/ResolutionLocalComparison.lean}
\CatalogueNote{Scope}{arbitrary u}
\CatalogueNote{Manuscript context}{cs.thm.equivalence}
\CatalogueNote{Construction and reuse}{Project-specific comparison of exact local and Reedy weak classes; derived mapping-space and two-out-of-three results are supporting inputs.}
\end{resultblock}

\begin{resultblock}
\begin{proposition}[Existence and exact classes of the resolution localization]
\label{catalogue:A7}
For the literal small family $S_h$ of free dagger collapse maps, the exact local weak-equivalence class admits a small solution set relative to the Reedy generating cofibrations.
Smith recognition therefore constructs a model structure on simplicial dagger diagrams whose cofibrations are exactly the Reedy cofibrations and whose weak equivalences are exactly the $S_h$-local equivalences.
Pushouts along cofibrations preserve local equivalences, and the localized model is locally presentable with two small generating families satisfying the small object argument.
Hence the specified resolution localization is left proper and combinatorial.
\end{proposition}

\Source{DaggerModels.PartII.resolutionLocalSolutionSet;
DaggerModels.PartII.resolutionLocalizationSmithData;
DaggerModels.PartII.resolutionLocalizedModelCategory;
DaggerModels.PartII.resolutionLocalized_cofibrations;
DaggerModels.PartII.resolutionLocalized_weakEquivalences;
DaggerModels.PartII.resolutionLocalized_leftProper;
DaggerModels.PartII.resolutionLocalized_combinatorial}{lean/DaggerModelsPartII/ResolutionLocalSolutionSet.lean;
lean/DaggerModelsPartII/ResolutionLocalizationSmith.lean;
lean/DaggerModelsPartII/ResolutionModel.lean}
\CatalogueNote{Scope}{arbitrary u}
\CatalogueNote{Manuscript context}{cs.def.Sh}
\CatalogueNote{Construction and reuse}{New Part II solution-set and exact-class assembly using the project's generic Smith machinery and the Part I dagger Joyal coefficient model.}
\end{resultblock}

\begin{resultblock}
\begin{proposition}[Constant diagrams and evaluation form a Quillen equivalence]
\label{catalogue:A8}
Equip $\sSet^{\dag}$ with the Part I dagger Joyal model and $\sSpace^{\dag\vee}$ with the specified $S_h$-localized Reedy model.
The adjunction $c\dashv\ev_0$ between constant diagrams and degree-zero evaluation is a Quillen equivalence.
Its proof uses the exact Reedy cofibrations and local weak equivalences of the constructed model, the characterization of its fibrant objects, and the comparison identifying local equivalences into a local target with weak equivalences of the adjunct map in degree zero.
\end{proposition}

\Source{DaggerModels.PartII.constantEvaluation_localEquivalence_iff;
DaggerModels.PartII.constantEvaluation_isQuillenEquivalence_of_exactClasses;
DaggerModels.PartII.resolutionLocalized_constantEvaluation}{lean/DaggerModelsPartII/ResolutionLocalComparison.lean;
lean/DaggerModelsPartII/ResolutionLocalizedModel.lean;
lean/DaggerModelsPartII/ResolutionModel.lean}
\CatalogueNote{Scope}{arbitrary u}
\CatalogueNote{Manuscript context}{cs.thm.equivalence}
\CatalogueNote{Construction and reuse}{Project-specific Quillen-equivalence criterion and application to the concrete localized model, relying on the Part I coefficient structure.}
\end{resultblock}
